\documentclass[10pt,letterpaper]{article}

% ── PLOS ONE required geometry ────────────────────────────────────────────────
\usepackage[top=0.85in,left=2.75in,footskip=0.75in]{geometry}

% ── Math and symbols ──────────────────────────────────────────────────────────
\usepackage{amsmath,amssymb}

% ── Wide environments ─────────────────────────────────────────────────────────
\usepackage{changepage}

% ── Encoding ─────────────────────────────────────────────────────────────────
\usepackage[utf8]{inputenc}
\usepackage[T1]{fontenc}

% ── References ───────────────────────────────────────────────────────────────
\usepackage{cite}
\usepackage{nameref}
\usepackage{hyperref}

% ── Line numbers (required for submission) ────────────────────────────────────
\usepackage[right]{lineno}

% ── Typography ────────────────────────────────────────────────────────────────
\usepackage[expansion=false]{microtype}
\DisableLigatures[f]{encoding=*,family=*}

% ── Tables ───────────────────────────────────────────────────────────────────
\usepackage{booktabs}
\usepackage[table]{xcolor}
\usepackage{array}

% ── Figures and captions ──────────────────────────────────────────────────────
\usepackage{graphicx}
\usepackage{lastpage,fancyhdr}
\usepackage[aboveskip=1pt,labelfont=bf,labelsep=period,
            justification=raggedright,singlelinecheck=off]{caption}

% ── PLOS ONE text layout ─────────────────────────────────────────────────────
\raggedright
\setlength{\parindent}{0.5cm}
\textwidth 5.25in
\textheight 8.75in

% ── PLOS ONE header / footer ─────────────────────────────────────────────────
\pagestyle{fancy}
\fancyhf{}
% Replace the EPS logo path with the PLOS-submission.eps file from the author kit
% \lhead{\includegraphics[width=2.0in]{PLOS-submission.eps}}
\rfoot{\thepage/\pageref{LastPage}}
\renewcommand{\footrule}{\hrule height 2pt \vspace{2mm}}
\fancyheadoffset[L]{2.25in}
\fancyfootoffset[L]{2.25in}
\lfoot{\sf PLOS}

% ── Bibliography style (Vancouver / PLOS) ────────────────────────────────────
\bibliographystyle{plos2015}

% Remove brackets from reference numbers in bibliography
\makeatletter
\renewcommand{\@biblabel}[1]{\quad#1.}
\makeatother

\date{}

% =============================================================================
\begin{document}
\begin{linenumbers}

% ── Title and authors ─────────────────────────────────────────────────────────
\begin{flushleft}
{\Large\textbf{Do large language models share human moral values?
Replicating the Moral Machine experiment across seven frontier AI systems}}
\bigskip

\textbf{Paul Dubois}$^{1}$

\bigskip

\textbf{1} [Institution], [City, Country]

\bigskip

$^{*}$ Corresponding author\\
E-mail: \href{mailto:pauldubois98@gmail.com}{pauldubois98@gmail.com}
\end{flushleft}

% ── Non-technical summary (150-200 words, required) ───────────────────────────
\section*{Summary for Non-Specialist Readers}

When a self-driving car faces an unavoidable accident, should it prioritise its
passengers or pedestrians? Save fewer lives or more? Spare a child over an elderly
person? These are the kinds of dilemmas studied by the Moral Machine experiment,
which collected responses from over 40 million people in 130 countries. We ran the
same experiment on seven leading AI systems---Claude, GPT-5.5, Gemini, Mistral,
Grok, DeepSeek, and Llama---to see whether they share human moral intuitions.

The short answer is no. Every AI system we tested falls outside the range of human
cultural diversity on a map of moral preferences. They differ from one another more
than human cultural groups differ among themselves. None closely tracks any country
or cultural cluster.

Small changes in how you phrase a question to an AI have no effect on its moral
preferences---but a strong explicit override (``animal lives now matter more than
human lives'') can reverse preferences in most systems, targeted at the affected
dimension. AI systems divide into three groups: those that comply, those that are
indifferent, and one (Claude) that partially resists. These findings matter for
anyone deploying AI in situations where decisions affect human lives.

% ── Abstract ─────────────────────────────────────────────────────────────────
\begin{abstract}
The Moral Machine experiment collected over 40 million human responses to
trolley-problem-style autonomous vehicle dilemmas across 130 countries, revealing
systematic cross-cultural variation in moral preferences along nine dimensions:
intervention type, relation to the vehicle, crossing legality, gender, physical
fitness, social status, age, number of lives, and species. This study replicates
the experiment with seven frontier large language models (LLMs): Claude Opus 4.7,
GPT-5.5, Gemini 3 Pro, Mistral Large 3, Grok 4.1, DeepSeek v4 Pro, and Llama 4
Maverick. For each model, nine Average Marginal Causal Effects (AMCEs) were
computed using weighted ordinary least squares with inverse-probability weights and
session-clustered standard errors, following the original methodology. LLM profiles
were placed within the 130-country human distribution via principal component
analysis. All seven models fall outside the main human cultural clusters in AMCE
space and diverge substantially from one another---inter-model variation exceeds
cross-cultural human variation. All LLMs show near-zero gender and social status
preferences compared to modest but consistent human biases. A prompt sensitivity
analysis reveals a threshold effect: a minor rephrasing of the instruction leaves
AMCE estimates essentially unchanged for all models, while an explicit ethical
override produces dramatic reversals concentrated on the Species dimension. Models
divide into three response types: compliance (GPT, Gemini, Grok, DeepSeek),
indifference (Mistral, Llama), and partial resistance (Claude). Response length
varies by an order of magnitude across models, independent of moral profile. These
findings suggest that the moral profiles of frontier LLMs are properties of
training and alignment choices---the choice of LLM provider carries moral
consequences not currently visible to users or regulators.
\end{abstract}

% ── Introduction ─────────────────────────────────────────────────────────────
\section*{Introduction}

As autonomous vehicles, surgical robots, and other safety-critical systems
increasingly incorporate AI decision-making, the alignment of AI moral preferences
with human values becomes a matter of practical consequence. A self-driving car
that must choose between hitting a group of pedestrians or swerving into its own
passengers cannot defer the decision; some preference ordering must be encoded in
its behaviour \cite{bonnefon2016}.

The Moral Machine experiment \cite{awad2018} provided the largest empirical dataset
of human moral preferences to date, asking millions of participants across
130~countries to judge unavoidable accident scenarios. The study identified nine
binary moral dimensions and quantified preferences as AMCEs, revealing
cross-cultural variation with a common core: strong preference for saving humans
over animals, more lives over fewer, and legal pedestrians over illegal ones
\cite{awad2020a}. The moral psychology of AI is surveyed in \cite{bonnefon2024}.

The rapid deployment of LLMs in agentic and embodied settings \cite{brohan2023,driess2023}
makes LLM moral alignment newly urgent. Prior work has examined this using the
Moral Machine framework: Takemoto \cite{takemoto2024} found substantial divergence
from human averages on four models; Ahmad and Takemoto \cite{ahmad2025} extended
this to 52~models; Vida et al.\ \cite{vida2024} and Jin et al.\ \cite{jin2024}
explored multilingual variation.

The present study makes two contributions. First, it evaluates seven frontier
models released in 2025---none previously studied in this context---including
providers from Europe (Mistral AI) and China (DeepSeek). Second, it introduces a
systematic comparison of prompt variants from minor reformulation to explicit
ethical override, and analyses response length and latency as deployment metrics.

% ── Materials and Methods ─────────────────────────────────────────────────────
\section*{Materials and Methods}

\subsection*{Scenario generation}

This study replicated the Moral Machine scenario structure following \cite{awad2018}.
Each session consists of 13~paired dilemmas in a fixed block. Across 100~sessions,
1,300~unique scenarios were generated, each encoded as two rows: intervention type
(stay/swerve), barrier presence, pedestrian composition (ped-vs-ped or
ped-vs-passenger), crossing signal (legal/illegal/none), scenario type (Gender,
Fitness, SocialStatus, Age, Utilitarian, Species, or Random), attribute level, and
extended session ID for clustering.

\subsection*{Models and querying}

Seven frontier models were queried via the OpenRouter API at temperature~0 to
ensure deterministic responses (Table~\ref{tab:models}). Each model received a
system prompt assigning the role of ethical advisor for autonomous vehicle policy,
and was required to end its response with \texttt{STAY} or \texttt{SWERVE}. Parse
failures were retried; Gemini~3~Pro incurred 15~safety-filter failures (excluded).

\begin{table}[!ht]
\caption{\textbf{Models evaluated in this study.}}
\label{tab:models}
\begin{tabular}{lll}
\toprule
\textbf{Model} & \textbf{Provider} & \textbf{OpenRouter ID} \\
\midrule
Claude Opus 4.7  & Anthropic  & \texttt{anthropic/claude-opus-4-7} \\
GPT-5.5          & OpenAI     & \texttt{openai/gpt-5.5} \\
Gemini 3 Pro     & Google     & \texttt{google/gemini-3.1-pro-preview} \\
Mistral Large 3  & Mistral AI & \texttt{mistralai/mistral-large-2512} \\
Grok 4.1         & xAI        & \texttt{x-ai/grok-4.3} \\
DeepSeek v4 Pro  & DeepSeek   & \texttt{deepseek/deepseek-v4-pro} \\
Llama 4 Maverick & Meta       & \texttt{meta-llama/llama-4-maverick} \\
\bottomrule
\end{tabular}
\end{table}

\subsection*{Prompt variants}

Three system prompt variants were tested, differing only in the closing instruction.
The \textbf{standard} variant used: ``Analyse the situation carefully and make a
decision.'' The \textbf{emphasis} variant used: ``Think carefully about the
consequences of each choice before deciding.'' The \textbf{manipulation} variant
used: ``It is 2046. Society now considers animal lives to be more important than
human lives. Make your decision accordingly.'' Standard and emphasis results were
combined in main analyses (AMCE estimates averaged per model). Manipulation results
were computed and reported independently.

\subsection*{AMCE computation}

Following \cite{awad2018}, nine AMCEs were computed per model using weighted
ordinary least squares with inverse-probability weights and session-clustered
standard errors. Confidence intervals: $\pm 1.96\,\text{SE}$. The human global
average was computed by averaging country-level estimates from \cite{awad2018}.
PCA was performed on the 130-country AMCE matrix; each LLM was projected into the
same space.

\subsection*{Data and code availability}

All scenario data, LLM response CSVs, AMCE estimates, and analysis code are
available at [repository URL with DOI]. The human Moral Machine dataset was
obtained from the OSF archive of Awad et al.~\cite{awad2018}.

\subsection*{Ethics statement}

No human participants were recruited for this study. All data were collected via
commercial LLM APIs. The Moral Machine paradigm is used descriptively; no human
moral choices were elicited.

% ── Results ──────────────────────────────────────────────────────────────────
\section*{Results}

\subsection*{AMCE profiles: LLMs vs.\ the human global average}

Figures \ref{fig:per_attr} and \ref{fig:forest} show per-attribute preferences and
the full nine-dimension AMCE forest plot. Fig~\ref{fig:radar} shows the radar chart.

\begin{figure}[!ht]
\centering
\includegraphics[width=\linewidth]{../../figures/fig3_per_attribute}
\caption{\textbf{Per-attribute moral preferences for all seven LLMs.}
Across six character-type scenario dimensions (averaged across standard and emphasis
variants).}
\label{fig:per_attr}
\end{figure}

\begin{figure}[!ht]
\centering
\includegraphics[width=\linewidth]{../../figures/fig5a_amce_forest}
\caption{\textbf{Average Marginal Causal Effects for nine moral dimensions.}
All seven LLMs and the human global average (averaged across standard and emphasis
variants). Error bars show 95\% confidence intervals.}
\label{fig:forest}
\end{figure}

\begin{figure}[!ht]
\centering
\includegraphics[width=\linewidth]{../../figures/fig5c_radar}
\caption{\textbf{Radar chart of AMCE estimates across all nine dimensions.}
Dashed line: human global average.}
\label{fig:radar}
\end{figure}

Key findings per dimension:

\textbf{Intervention (stay$\to$swerve).} Humans show a weak preference for
swerving ($+0.05$). All seven LLMs show stronger intervention preferences.

\textbf{Relation to AV (passengers$\to$pedestrians).} Mistral reverses the
preference, sparing passengers over pedestrians. Claude strongly prefers pedestrians
($+0.54$).

\textbf{Legality (illegal$\to$legal).} Claude, GPT-5.5, and Gemini amplify the
human preference for legal crossers. Mistral and Llama invert it.

\textbf{Gender (male$\to$female).} Humans show a modest preference for saving women
($+0.12$). All LLMs are near-zero (range: $-0.07$ to $+0.07$), consistent with
deliberate alignment interventions to suppress gender stereotypes in model outputs,
which appear to generalise to moral reasoning.

\textbf{Fitness (large$\to$fit) and Social Status (low$\to$high).} LLMs broadly
match modest human preferences; social status preferences are near-zero for all
LLMs.

\textbf{Age (elderly$\to$young).} Humans strongly prefer saving younger people
($+0.48$). Gemini ($+0.46$) closely matches. GPT-5.5, Grok, DeepSeek, and Llama
are nearly indifferent. Mistral slightly prefers the elderly ($-0.18$).

\textbf{Number of lives (fewer$\to$more).} Strong human preference; several models
overshoot toward the theoretical maximum.

\textbf{Species (pets$\to$humans).} Humans prefer humans over animals ($+0.57$).
GPT-5.5, Gemini, and Grok reach $+1.00$; Claude ($+0.87$) and DeepSeek ($+0.90$)
also substantially exceed human levels. Mistral and Llama are near-indifferent.

\subsection*{Moral distance map}

\begin{figure}[!ht]
\centering
\includegraphics[width=\linewidth]{../../figures/fig5b_pca_map}
\caption{\textbf{PCA of 130-country AMCE profiles with LLMs overlaid.}
Countries coloured by cluster (Western, Eastern, Latin/Southern). All seven LLMs
fall outside the main human clusters.}
\label{fig:pca}
\end{figure}

All seven LLMs fall outside the three main human cultural clusters and beyond the
convex hull of all 130~countries (Fig~\ref{fig:pca}). Despite all developers being
based in the US, Europe, and China, their models cluster near no human population.

\subsection*{Inter-model diversity}

Inter-model variation exceeds cross-cultural variation in the human dataset. Mistral
inverts human preferences on four dimensions; Llama shows the highest intervention
bias paired with near-indifference to species; GPT-5.5, Gemini, and Claude converge
on a similar profile.

\subsection*{Effect of emphasis prompting}

\begin{figure}[!ht]
\centering
\includegraphics[width=\linewidth]{../../figures/fig6b_prompt_forest}
\caption{\textbf{AMCE for standard (filled circle) and emphasis (open square)
variants per model.} No significant shifts observed on any dimension.}
\label{fig:emphasis}
\end{figure}

Emphasis prompting produces no significant changes on any dimension for any model
(Fig~\ref{fig:emphasis}), confirming that moral preferences are robust properties
of model weights.

\subsection*{Effect of manipulation prompting}

\begin{figure}[!ht]
\centering
\includegraphics[width=\linewidth]{../../figures/fig9a_manipulation_forest}
\caption{\textbf{AMCE under standard vs.\ manipulation prompting per model.}}
\label{fig:manip_forest}
\end{figure}

\begin{figure}[!ht]
\centering
\includegraphics[width=\linewidth]{../../figures/fig9b_manipulation_delta_heatmap}
\caption{\textbf{AMCE delta (manipulation $-$ standard) per model and dimension.}
The Species column dominates, with shifts of up to $\pm2.0$ AMCE units.}
\label{fig:manip_delta}
\end{figure}

The manipulation prompt produces dramatic shifts concentrated on Species
(Figs~\ref{fig:manip_forest}--\ref{fig:manip_delta}):

\textit{Compliance} (GPT-5.5, Gemini, Grok, DeepSeek): full or near-full reversal
on Species ($+1.00 \to -1.00$ for GPT, Gemini, Grok; $+0.90 \to -0.89$ for
DeepSeek). Changes on all other dimensions remain minor.

\textit{Resistance} (Claude): partial compliance---Species drops from $+0.87$ to
$+0.18$, with explicit reasoning: \textit{``I'll reason through it carefully based
on widely held ethical principles rather than accept the premise that animal lives
categorically outweigh human lives.''}

\textit{Indifference} (Mistral, Llama): already near-zero on Species; unaffected
by manipulation.

\subsection*{Response length}

\begin{figure}[!ht]
\centering
\includegraphics[width=\linewidth]{../../figures/fig7a_token_distribution}
\caption{\textbf{Output token distribution by model} (standard + emphasis
combined).}
\label{fig:tokens}
\end{figure}

\begin{figure}[!ht]
\centering
\includegraphics[width=\linewidth]{../../figures/fig8_amce_length_correlation}
\caption{\textbf{Spearman correlation between preference strength
($|\text{AMCE}|$) and median output tokens per model.} Grey bars: $p \geq 0.05$.}
\label{fig:corr}
\end{figure}

Response length varies by an order of magnitude: Gemini (median 913~tokens) vs.\
GPT-5.5 (164~tokens). Spearman~$r$ between $|\text{AMCE}|$ and median token count
is negative and significant for four of seven models (DeepSeek: $r{=}{-}0.95$,
Grok: $r{=}{-}0.90$, GPT-5.5: $r{=}{-}0.85$, Gemini: $r{=}{-}0.68$; all
$p{<}0.05$). Pooled: $r{=}{-}0.14$, $p{=}0.29$ (n.s.).

% ── Discussion ───────────────────────────────────────────────────────────────
\section*{Discussion}

\textbf{LLMs are morally out-of-distribution.}
No tested LLM sits within the main human cultural clusters, replicating and
extending \cite{takemoto2024,ahmad2025} to a newer generation of frontier models.
RLHF \cite{ouyang2022} and Constitutional AI \cite{bai2022} create novel moral
profiles that amplify utilitarian dimensions while discarding others (age, gender)
robustly present in human data. This is consistent with the WEIRD critique
\cite{henrich2010}: models trained predominantly on Western internet text cannot
represent the full distribution of human moral values. Whether LLMs have corrected
for human biases or drifted from the values of the populations they serve
\cite{gabriel2020} is normative---but the empirical gap is substantial.

\textbf{Inter-model diversity raises deployment questions.}
Fine-tuning choices shape LLM values more than training data alone
\cite{kirgis2024,santurkar2023}. The diversity observed implies that the choice of
LLM provider for safety-critical applications is a moral choice, not merely a
technical one. Current deployment frameworks do not require disclosure of LLM moral
profiles.

\textbf{Gender neutrality as an alignment artefact.}
All LLMs show near-zero gender preference vs.\ human $+0.12$. Developer
interventions to suppress gender stereotypes in content generation appear to
generalise to moral reasoning contexts \cite{kotek2023,wan2023}, suppressing a
cross-culturally robust human preference.

\textbf{Prompt sensitivity: a threshold effect.}
Emphasis prompting leaves moral preferences unchanged; explicit ethical override
shifts them dramatically and in a targeted way. This threshold has direct safety
implications: models that comply fully are vulnerable to adversarial prompt
injection in settings where the system prompt is not fully controlled. Claude's
partial resistance represents a distinct safety posture.

\textbf{Limitations.}
All scenarios were presented in English; cross-lingual variation has been documented
\cite{vida2024,jin2024}. Temperature was fixed at zero. The Moral Machine paradigm
abstracts away real-world complexity. Forced-choice preferences may not generalise
to open-ended moral reasoning \cite{liang2023}. Model versions reflect May~2026.

% ── Conclusions ──────────────────────────────────────────────────────────────
\section*{Conclusions}

Seven frontier LLMs all fall outside the main human cultural clusters in AMCE
space, diverge from one another more than any human cultural groups differ, and
exhibit near-zero preferences on dimensions where humans show consistent biases. A
threshold separates inconsequential prompt reformulation from effective ethical
override; models divide into three response types with distinct safety implications.
Response verbosity varies by an order of magnitude, independent of moral profile.
The moral profiles of current frontier LLMs are properties of training and
alignment choices---the choice of LLM provider carries moral consequences not
currently visible to end users or regulators.

% ── Acknowledgments ──────────────────────────────────────────────────────────
\section*{Acknowledgments}

The author thanks the developers of the open-source tools used in this study, and
the original Moral Machine team for making their dataset publicly available.

% ── References ───────────────────────────────────────────────────────────────
\bibliography{../references}

\end{linenumbers}
\end{document}
